% Partie II condensée (FR). Une seule figure : le compte-rendu de sortie (signature).
% Les registres article/forum sont évoqués en prose (contrainte 15 pages).

\noindent
\begin{center}
\begin{tikzpicture}
\node[draw=ThesisNeutral, fill=ThesisPaper, rounded corners=3pt, line width=0.5pt,
      inner sep=10pt] {%
\begin{tabular}{@{}l !{\hspace{6pt}\textcolor{ThesisNeutral}{\vrule width 0.3pt}\hspace{6pt}} l@{}}
\begin{minipage}[t]{0.40\textwidth}
\scriptsize\color{ThesisInk}
{\tiny\sffamily\color{ThesisNeutral} FRANÇAIS}\\[4pt]
{\bfseries CENTRE HOSPITALIER DE SAINT-AUBIN}\\ Dossier n\textsuperscript{o} 4471902\\[1pt]
Cardiologie, Soins Intensifs\\ Séjour : 03/11--08/11/2024\\[4pt]
{\color{ThesisNeutral}\rule{\linewidth}{0.3pt}}\\[4pt]
\textbf{Patient :} CARON Julien, né le 14/03/1962 (62 ans). \textbf{Dest. :} D\textsuperscript{r} S.~Lemaire.\\[5pt]
\textbf{Motif.} Douleur thoracique constrictive rétrosternale, irradiation bras G, depuis 2h.\\[5pt]
\textbf{Antécédents.} HTA. DT2 sous metformine. Tabac 30 PA. Dyslipidémie.\\[5pt]
\textbf{Examen.} TA 148/92, FC 98, SpO\textsubscript{2} 96\% AA, EVA 7/10. Pas de signe d'IVG.\\[5pt]
\textbf{Paraclinique.} ECG : sus-décalage ST V1--V4. Tropo T 4520 ng/L (N $<$14). Coro : occlusion IVA proximale, angioplastie + stent actif, TIMI 3. ETT : FEVG 40\%.\\[5pt]
\textbf{Au total.} SCA ST+ antérieur (I21.0).\\[5pt]
\textbf{Ttt sortie.} Kardégic 75, ticagrélor 90x2, atorvastatine 80, bisoprolol 2,5, ramipril 5.\\[5pt]
\textbf{CAT.} Réadaptation cardiaque. Consultation cardio à 1 mois.\hfill D\textsuperscript{r} C.~Després
\end{minipage}
&
\begin{minipage}[t]{0.40\textwidth}
\scriptsize\color{ThesisNeutral}
{\tiny\sffamily\color{ThesisNeutral} ENGLISH · TRANSLATION}\\[4pt]
{\bfseries SAINT-AUBIN HOSPITAL}\\ Record no. 4471902\\[1pt]
Cardiology, Intensive Care\\ Stay: 03/11--08/11/2024\\[4pt]
{\color{ThesisNeutral}\rule{\linewidth}{0.3pt}}\\[4pt]
\textbf{Patient:} CARON Julien, b.\ 14/03/1962 (62 y). \textbf{To:} Dr S.~Lemaire.\\[5pt]
\textbf{Reason.} Constrictive retrosternal chest pain, radiating to L arm, for 2h.\\[5pt]
\textbf{History.} HTN. T2DM on metformin. Smoking 30 PY. Dyslipidaemia.\\[5pt]
\textbf{Exam.} BP 148/92, HR 98, SpO\textsubscript{2} 96\% RA, pain 7/10. No sign of LVF.\\[5pt]
\textbf{Workup.} ECG: ST elevation V1--V4. Trop T 4520 ng/L (N $<$14). Angiography: proximal LAD occlusion, angioplasty + drug-eluting stent, TIMI 3. Echo: LVEF 40\%.\\[5pt]
\textbf{Summary.} Anterior STEMI (I21.0).\\[5pt]
\textbf{Discharge tx.} Aspirin 75, ticagrelor 90x2, atorvastatin 80, bisoprolol 2.5, ramipril 5.\\[5pt]
\textbf{Plan.} Cardiac rehabilitation. Cardiology review at 1 month.\hfill Dr C.~Després
\end{minipage}
\end{tabular}%
};
\end{tikzpicture}
\captionof{figure}{Un compte-rendu de sortie pour un infarctus du myocarde aigu ; original français, traduction anglaise.}
\label{fig:bt-clinical-note}
\end{center}
\medskip

\noindent La \Cref{fig:bt-clinical-note} est une lettre d'un cardiologue à un confrère. Pour un
lecteur français ordinaire, elle est très difficile à lire : peu de verbes, des
phrases télégraphiques, une information dense simplement séparée par des virgules
comme dans un tableau. Ce n'est pas de la négligence, mais un style de communication
spécialement ajusté à la communauté à laquelle il s'adresse. Le linguiste Zellig
Harris a appelé une telle forme spécialisée et régie par ses propres règles un
\textit{sous-langage}~\citep{harris_mathematical_1968}, et le compte-rendu clinique
en est le cas le mieux documenté~\citep{sager_nlip_1981}. C'est avant tout une forme
de compression : \citet{anderson_compression_1975} a montré que l'abrègement clinique
obéit à un petit ensemble de règles de suppression, ne retirant que ce que le lecteur
peut rétablir. Dans la \Cref{fig:bt-clinical-note}, \emph{Pas de signe d'IVG} tient
pour \emph{il n'y a pas de signe d'insuffisance ventriculaire gauche} ; ce qui
disparaît est la part que le lecteur expert connaît déjà~\citep{stalnaker_common_2002}.
Le compte-rendu est aussi performatif : une ordonnance permet au patient d'obtenir le
médicament, un ordre de transfert enclenche le transfert~\citep{austin_words_1962}.

Tout cela rend le compte-rendu précieux à étudier, mais aussi difficile à obtenir. Il
quitte rarement l'hôpital qui l'a produit, car il est rempli de données personnelles,
et retirer les identifiants évidents ne suffit pas : un compte-rendu pseudonymisé
n'est pas anonyme. Pour le français, le problème est pire. Il n'existe pas de corpus
partageable de vrais dossiers hospitaliers français~\citep{neveol_clinical_2018}. Les
données publiques sont une alternative : cas cliniques publiés, notices de
médicaments, résumés d'articles. Mais un cas clinique est écrit pour enseigner ; il se
lit comme un récit, en phrases complètes et au passé. Le même infarctus, raconté dans
un article de recherche ou sur un forum de patients (\Cref{fig:bt-web}), est abondant
et public, mais il parle une autre langue.

\begin{center}
\begin{tikzpicture}
\node[draw=ThesisNeutral, fill=ThesisPaper, rounded corners=3pt, line width=0.5pt,
      inner sep=10pt] {%
\begin{tabular}{@{}l !{\hspace{6pt}\textcolor{ThesisNeutral}{\vrule width 0.3pt}\hspace{6pt}} l@{}}
\begin{minipage}[t]{0.40\textwidth}
\footnotesize\color{ThesisInk}
{\tiny\sffamily\color{ThesisNeutral} FRANÇAIS}\\[5pt]
\textit{Bonjour, mon père a fait une crise cardiaque la semaine dernière. Il a eu
une grosse douleur dans la poitrine qui descendait dans le bras gauche, il était
tout pâle et en sueur, on a appelé le 15 tout de suite. Est-ce que quelqu'un sait
combien de temps dure la convalescence après un infarctus ? Merci d'avance.}
\end{minipage}
&
\begin{minipage}[t]{0.40\textwidth}
\footnotesize\color{ThesisNeutral}
{\tiny\sffamily\color{ThesisNeutral} ENGLISH · TRANSLATION}\\[5pt]
\textit{Hello, my father had a heart attack last week. He had a strong pain in the
chest that went down into the left arm, he was all pale and sweating, we called
emergency services right away. Does anyone know how long recovery takes after a
heart attack? Thanks in advance.}
\end{minipage}
\end{tabular}%
};
\end{tikzpicture}
\captionof{figure}{Un forum de patients à propos de l'infarctus du myocarde ; original français, traduction anglaise.}
\label{fig:bt-web}
\end{center}
\medskip

Un texte biomédical n'est pas une langue unique. La même connaissance médicale change
selon qui écrit, qui lit et ce que le texte est censé faire. Comptes-rendus cliniques,
articles de recherche et textes destinés aux patients offrent donc des vues
différentes de la médecine, et seules les deux dernières sont largement disponibles.
Le défi pour la suite de cette thèse est d'utiliser ce texte public sans perdre de vue
la langue clinique dont nous avons finalement besoin.
